\section{Filters and ultrafilters}
\label{sec:types-ultraproducts-04}

\begin{definition}[Filter]
Let \(I\) be a nonempty set.  A \emph{filter} on \(I\) is a set
\(\mathcal F\subseteq\Pow(I)\) such that:
\begin{enumerate}
  \item \(I\in\mathcal F\) and \(\varnothing\notin\mathcal F\);
  \item if \(X,Y\in\mathcal F\), then \(X\cap Y\in\mathcal F\); and
  \item if \(X\in\mathcal F\) and \(X\subseteq Y\subseteq I\), then
        \(Y\in\mathcal F\).
\end{enumerate}
An \emph{ultrafilter} is a filter \(\mathcal U\) such that, for every
\(X\subseteq I\), exactly one of \(X\) and \(I\setminus X\) belongs to
\(\mathcal U\).  The requirement that exactly one of \(X\) and
\(I\setminus X\) belong to \(\mathcal U\) is cited below as the
\emph{ultrafilter alternative}.
\end{definition}

An ultrafilter is \emph{principal} iff it has the form
\(\set{X\subseteq I:i_0\in X}\) for some \(i_0\in I\).  A subset
\(X\subseteq I\) is \emph{cofinite} iff \(I\setminus X\) is finite.

\begin{lemma}[Existence of nonprincipal ultrafilters]
\label{lem:nonprincipal-ultrafilter}
Assume the Ultrafilter Lemma in the ambient set theory.  Let \(I\) be an
infinite set.
\begin{enumerate}
  \item The cofinite subsets of \(I\) form a filter on \(I\), the
        \emph{Fr\'echet filter} on \(I\).
  \item Some ultrafilter on \(I\) contains the Fr\'echet filter, and every
        ultrafilter containing it is nonprincipal.
  \item An ultrafilter on \(I\) is nonprincipal iff it contains every
        cofinite subset of \(I\).
\end{enumerate}
In particular, there is a nonprincipal ultrafilter on \(\N\).
\textit{Status: complete proof.}
\end{lemma}

\begin{proof}
For clause 1, check the three filter conditions.  \(I\setminus
I=\varnothing\) is finite, so \(I\) is cofinite, while
\(I\setminus\varnothing=I\) is infinite, so \(\varnothing\) is not.  If
\(X\) and \(Y\) are cofinite, then
\(I\setminus(X\cap Y)=(I\setminus X)\cup(I\setminus Y)\) is a union of two
finite sets, hence finite, so \(X\cap Y\) is cofinite.  If \(X\) is
cofinite and \(X\subseteq Y\subseteq I\), then
\(I\setminus Y\subseteq I\setminus X\) is finite, so \(Y\) is cofinite.

For clause 2, the Fr\'echet filter satisfies the two displayed conditions
of the Ultrafilter Lemma (\cref{def:ultrafilter-lemma}); these are exactly
the filter conditions just verified.  The assumed principle therefore
supplies an ultrafilter containing the Fr\'echet filter.  Now let
\(\mathcal U\) be any ultrafilter containing the Fr\'echet filter and let
\(i_0\in I\).  The set \(I\setminus\set{i_0}\) is cofinite, so it belongs
to \(\mathcal U\); it omits \(i_0\), so
\(\mathcal U\ne\set{X\subseteq I:i_0\in X}\).  Since \(i_0\) was
arbitrary, \(\mathcal U\) is nonprincipal.

For clause 3, right to left, an ultrafilter containing every cofinite
subset contains the Fr\'echet filter and is nonprincipal by clause 2.
Left to right, argue contrapositively in three steps that an ultrafilter
\(\mathcal U\) containing some finite set is principal.  First, if
\(X\cup Y\in\mathcal U\), then \(X\in\mathcal U\) or \(Y\in\mathcal U\):
otherwise the ultrafilter alternative would put \(I\setminus X\) and
\(I\setminus Y\) in \(\mathcal U\), hence also
\((I\setminus X)\cap(I\setminus Y)=I\setminus(X\cup Y)\), and the
alternative would then exclude \(X\cup Y\), contrary to assumption.
Second, by induction on \(m\geq 1\): if \(F\in\mathcal U\) is finite with
\(\abs{F}=m\), then \(\set{i}\in\mathcal U\) for some \(i\in F\).  For
\(m=1\) the set \(F\) is itself a singleton.  For the step from \(m\) to
\(m+1\), choose \(i\in F\) and write
\(F=(F\setminus\set{i})\cup\set{i}\); the first step gives
\(\set{i}\in\mathcal U\), in which case the claim holds, or
\(F\setminus\set{i}\in\mathcal U\), a finite member of size \(m\) to which
the induction hypothesis applies.  Third, if \(\set{i_0}\in\mathcal U\),
then \(\mathcal U=\set{X\subseteq I:i_0\in X}\): every \(X\) with
\(i_0\in X\) includes \(\set{i_0}\) and belongs to \(\mathcal U\) by
upward closure, and no \(X\) with \(i_0\notin X\) belongs to
\(\mathcal U\), since otherwise \(X\cap\set{i_0}=\varnothing\) would lie
in \(\mathcal U\).  Contrapositively, a nonprincipal ultrafilter contains
no finite set, and the ultrafilter alternative then puts the complement of
every finite set, that is every cofinite set, into it.

The final claim is clause 2 applied to the infinite set \(\N\).
\end{proof}

